\documentclass[10pt,a4paper]{article}

% --- PACOTES BÁSICOS E TIPOGRAFIA ---
\usepackage[utf8]{inputenc}
\usepackage[T1]{fontenc}
\usepackage[margin=0.46in, top=0.40in, bottom=0.40in]{geometry}
\usepackage{titlesec}
\usepackage{enumitem}
\usepackage{hyperref}
\usepackage{xcolor}

% --- CORES E LINKS ---
\definecolor{primary}{RGB}{20, 24, 33}      % Quase preto para máxima legibilidade
\definecolor{accent}{RGB}{25, 75, 140}      % Azul sóbrio profissional
\definecolor{subtext}{RGB}{85, 95, 110}     % Cinza escuro para datas/locais

\hypersetup{
    colorlinks=true,
    linkcolor=accent,
    urlcolor=accent,
    pdfauthor={Samuel Amadeu Loureiro},
    pdftitle={Curriculum - Samuel Amadeu Loureiro - Machine Learning Engineering Intern}
}

% --- FORMATAÇÃO DE PÁGINA E ESPAÇAMENTOS ---
\pagestyle{empty}
\raggedbottom
\raggedright
\setlength{\tabcolsep}{0in}

% --- FORMATAÇÃO DE SEÇÕES ---
\titleformat{\section}
  {\bfseries\large\color{accent}}
  {}{0em}
  {}[\vspace{1.5pt}\color{accent}\titlerule]

\titlespacing*{\section}{0pt}{4.5pt}{2.5pt}

% --- LISTAS E BULLETS COMPACTOS ---
\setlist[itemize]{
    leftmargin=1.2em,
    label={\small\textbullet},
    itemsep=0.8pt,
    topsep=0.5pt,
    parsep=0pt,
    partopsep=0pt
}

\begin{document}

% ==============================================================================
% CABEÇALHO / HEADER
% ==============================================================================
\begin{center}
    {\Huge\bfseries\color{primary} Samuel Amadeu Loureiro}\\[2.5pt]
    \small\color{subtext}
    Brest/Plouzané, France
    \quad|\quad +33 07 43 52 84 87
    \quad|\quad \href{mailto:s25amadeu@enib.fr}{s25amadeu@enib.fr}
    \quad|\quad \href{linkedin.com/in/samuelamadeuloureiro}{LinkedIn}
    \quad|\quad \href{github.com/SamuelLoureiro}{GitHub}
\end{center}
\vspace{-6pt}

% ==============================================================================
% OBJETIVO / RESUMO PROFISSIONAL
% ==============================================================================
\section{Professional Summary}
\small\color{primary}
Computer Engineering student and Software Development enthusiast applying for the 2027 Software Dev Engineer Intern - Luxembourg role. Proficient in Python, C/C++, and distributed systems, with a proven record of building scalable cross-platform software and machine learning pipelines.

% ==============================================================================
% EDUCAÇÃO / FORMAÇÃO ACADÊMICA
% ==============================================================================
\section{Education}
\noindent
\textbf{\color{primary} Master in Computer Science: Intelligent and Autonomous Interactive Systems (SIIA)}, UBO, ENIB, ENSTA Bretagne, IMT Atlantique \hfill {\small\color{subtext} 09/2026 - 08/2027}
\newline {\footnotesize\color{subtext} Plouzané, France}
\begin{itemize}
    \item \small Relevant Coursework: Interactive Machine Learning and Deep Learning, Collective Intelligence and Multiagent Systems
\end{itemize}
\noindent
\textbf{\color{primary} Engineering Degree (Diplôme d'Ingénieur - Bac +5)}, École Nationale d'Ingénieurs de Brest (ENIB) \hfill {\small\color{subtext} 09/2025 - 08/2027}
\newline {\footnotesize\color{subtext} Plouzané, France}
\begin{itemize}
    \item \small Relevant Coursework: Network and Systems Communication, Interactive Applications Design, Microprocessors and Embedded Systems
\end{itemize}
\noindent
\textbf{\color{primary} Bachelor in Computer Engineering}, Universidade Federal de Itajubá (UNIFEI) \hfill {\small\color{subtext} 03/2022 - 06/2025}
\newline {\footnotesize\color{subtext} Itajubá, Brazil}
\begin{itemize}
    \item \small Relevant Coursework: Data Structures and Algorithms, Software Engineering, Databases, Algorithm Analysis
\end{itemize}

% ==============================================================================
% EXPERIÊNCIA PROFISSIONAL / ACADÊMICA
% ==============================================================================
\section{Work Experience}
\noindent
\textbf{\color{primary} Mobile Development \& Data Science Manager}, Asimov Jr. \hfill {\small\color{subtext} 12/2023 - 08/2025}
\newline {\footnotesize\color{subtext} MG, Brazil}
\newline {\footnotesize\textbf{Technologies:} \textit{Flutter / Dart, Python (Pandas, NumPy), Power BI, MongoDB, REST APIs}}
\begin{itemize}
    \item \small Architected cross-platform mobile compliance and tracking applications using Flutter, Dart, and MongoDB, handling REST APIs to serve commercial clients.
    \item \small Engineered automated data processing pipelines and business insight tools using Python (Pandas, NumPy) and Power BI as part of the newly structured Data Science department.
\end{itemize}
\noindent
\textbf{\color{primary} Embedded Systems Researcher (Team Professionalizing Project)}, ENIB \hfill {\small\color{subtext} 09/2025 - 12/2025}
\newline {\footnotesize\color{subtext} Plouzané, France}
\newline {\footnotesize\textbf{Technologies:} \textit{C, FreeRTOS, Raspberry Pi Pico (RP2350), SPI / UART / I2C, EasyEDA}}
\begin{itemize}
    \item \small Engineered real-time firmware in C using FreeRTOS for a dual Raspberry Pi Pico 2 data acquisition system, implementing deterministic master-slave SPI and UART protocols.
    \item \small Designed multilayer custom PCBs using EasyEDA, maintaining strict signal integrity and ground plane layouts.
\end{itemize}

% ==============================================================================
% PROJETOS EM DESTAQUE
% ==============================================================================
\section{Featured Projects}
\noindent
\textbf{\color{primary} TinyML: Connected Bovine Collar}
\hfill {\footnotesize\href{https://github.com/SamuelLoureiro}{[View Project / Code]}}
\newline {\footnotesize\textbf{Stack:} \textit{Arduino, Edge Impulse, C++, Machine Learning}}
\begin{itemize}
    \item \small Trained TinyML models using Edge Impulse to achieve 80\% accuracy in real-time classification of bovine behaviors from sensor movement data.
    \item \small Developed C++ firmware on Arduino devices for real-time data capture and on-device machine learning inference.
\end{itemize}
\noindent
\textbf{\color{primary} CAN Bus Sensor Network}
\hfill {\footnotesize\href{https://lnkd.in/es9Zi695}{[View Project / Code]}}
\newline {\footnotesize\textbf{Stack:} \textit{C / C++, STM32 Nucleo, CAN Bus, I2C, MPU9250}}
\begin{itemize}
    \item \small Developed C/C++ firmware within the STM32 ecosystem to centralize distributed multi-sensor data over a CAN bus architecture.
    \item \small Integrated an Invensense MPU9250 IMU via I2C to capture and encapsulate inertial data into structured CAN frames.
\end{itemize}
\noindent
\textbf{\color{primary} ROS 2 Mobile Robot}
\newline {\footnotesize\textbf{Stack:} \textit{ROS 2, C++ / Python, STM32, Raspberry Pi, Computer Vision}}
\begin{itemize}
    \item \small Implemented a distributed robotics network across STM32, Raspberry Pi, and Host PC communicating seamlessly via ROS 2.
    \item \small Engineered control interfaces and computer vision tracking algorithms supporting manual and autonomous operation modes.
\end{itemize}

% ==============================================================================
% HABILIDADES TÉCNICAS
% ==============================================================================
\section{Technical Skills}
\small
\noindent\textbf{Programming Languages:} Python, C / C++, Java, Dart, JavaScript, SQL\\[1.5pt]
\noindent\textbf{Frameworks \& Libraries:} PyTorch, TensorFlow, Scikit-Learn, Pandas, NumPy, OpenCV, Flutter\\[1.5pt]
\noindent\textbf{Databases \& Systems:} MongoDB, PostgreSQL, SQLite, ROS 2, FreeRTOS, Git / GitHub, Docker\\[1.5pt]

% ==============================================================================
% IDIOMAS E CERTIFICAÇÕES
% ==============================================================================
\section{Languages \& Certifications}
\small
\noindent\textbf{Languages:}
Portuguese (Native), English (Advanced (C1 - TOEIC)), French (Advanced)\\[1.5pt]

\end{document}